Let \(K\) be a compact subset of \(\mathcal R_{\mathrm{low}}\).  Choose a
nonnegative \(C^\infty\) function \(\psi\), supported on
\(\{z:\lVert z\rVert_\infty\leq1/2\}\), with \(\psi(0)=1\).  For a constant
\(c_*>0\) to be chosen, put
\[
H=n^{-1/(2\gamma+d)},\qquad
\delta=c_*H^\gamma,\qquad
\tau_1(x)=\delta\,\psi\!\left(\frac{x-x_0}{H}\right),\qquad
\tau_0(x)=0.
\]
Compactness of \(K\) and interiority of \(x_0\) give a common \(n_0\) for
which the support of the bump is contained in \([0,1]^d\).  The scaled-bump
calculation in the definition of the H\"older ball gives, uniformly over
\(K\),
\[
\lVert\tau_1\rVert_{\mathcal H^\gamma_d}\leq C_Kc_*.
\]
Furthermore \(\beta<\gamma\) on \(K\), with the gap bounded away from zero,
so
\[
\left\lVert\frac{\tau_1}{2}\right\rVert_{\mathcal H^\beta_d}
\leq C_Kc_*H^{\gamma-\beta}\leq C_Kc_*.
\]
Choose \(c_*\) small enough that these norms are at most \(L\) and all means
below lie in \([-1/2,1/2]\).

We now construct two full-data laws.  Under both, \(X\) is uniform on
\([0,1]^d\), \(A\mid X\) is Bernoulli with probability \(1/2\), and \(A\) is
conditionally independent of the potential outcomes.  Under \(P_j\), set
\[
\mu_{1,j}(x)=\frac{\tau_j(x)}2,\qquad
\mu_{0,j}(x)=-\frac{\tau_j(x)}2,\qquad j\in\{0,1\},
\]
and, conditionally on \(X=x\), let \(Y(a)\in\{-1,1\}\) have mean
\(\mu_{a,j}(x)\); the two potential outcomes may be taken conditionally
independent.  Finally define \(Y=AY(1)+(1-A)Y(0)\).  Thus consistency and
conditional exchangeability hold by construction.  The propensity is the
constant \(1/2\), so overlap and every propensity H\"older condition hold.
The preceding scaled-bump bounds give the outcome-regression and CATE
H\"older conditions.  The potential outcomes are bounded, and the uniform
density is normalized and satisfies \(f_-<1<f_+\).  Hence \(P_0,P_1\) both
belong to \(\mathcal P_{\mathrm{band},n}\).

It remains to bound their statistical distance without invoking an
unproved testing result.  Use the convention
\(H^2(Q_1,Q_0)=\int(\sqrt{q_1}-\sqrt{q_0})^2\).  Conditional on \(X=x\) and
\(A=a\), the outcome is a Rademacher variable whose mean changes from zero
to \(m_a(x)\), where \(|m_a(x)|=|\tau_1(x)|/2\).  Its chi-square divergence
from the mean-zero Rademacher law equals \(m_a(x)^2\), and squared Hellinger
distance is no larger than chi-square divergence.  Averaging over \(X,A\)
therefore yields
\[
H^2(P_{1,O},P_{0,O})\leq C\int\tau_1(x)^2\,dx
\leq C_K\delta^2H^d.
\]
Hellinger affinity tensorizes, so the identity
\[
H^2(P_{1,O}^{\otimes n},P_{0,O}^{\otimes n})
=2-2\left(1-\frac{H^2(P_{1,O},P_{0,O})}{2}\right)^n
\]
and \(1-(1-z)^n\leq nz\) imply
\[
H^2(P_{1,O}^{\otimes n},P_{0,O}^{\otimes n})
\leq C_Kn\delta^2H^d=C_Kc_*^2.
\]
Choose \(c_*\) still smaller so that the last Hellinger distance is at most
\(1/2\).  Cauchy--Schwarz applied to
\(|p-q|=|\sqrt p-\sqrt q|(\sqrt p+\sqrt q)\) gives total variation no larger
than Hellinger distance, hence at most \(1/2\).

For any estimator \(T_n\), the point targets are separated by
\(\tau_1(x_0)-\tau_0(x_0)=\delta\).  Pointwise, the triangle inequality gives
\(|T_n-\tau_0(x_0)|+|T_n-\tau_1(x_0)|\geq\delta\).  Integrating this inequality
against the minimum of the two \(n\)-sample densities yields
\[
\max_{j\in\{0,1\}}\mathbb E_{P_j}|T_n-\tau_j(x_0)|
\geq\frac{\delta}{2}\{1-\operatorname{TV}(P_{1,O}^{\otimes n},P_{0,O}^{\otimes n})\}
\geq\frac{\delta}{4}.
\]
Taking the infimum over \(T_n\) proves the result because
\(\delta=c_*n^{-\gamma/(2\gamma+d)}\).  All constants were chosen uniformly
over \(K\).